% =========================================================================
%  Capítulo 3 — Arquitetura
% =========================================================================

\chapter{Arquitetura}
\label{cap:arquitetura}

% --- INSTRUÇÕES (apagar antes da entrega) ---------------------
% Este capítulo é o que distingue um relatório bom de um
% relatório medíocre. Não basta dizer ``usámos Node.js e
% MongoDB'' — é preciso JUSTIFICAR cada decisão.
%
% Para cada decisão de arquitetura, respondam a três perguntas:
%   1. Que dados/lógica vivem aqui?
%   2. Porquê AQUI e não noutra camada?
%   3. Que alternativas foram consideradas e descartadas?
%
% Usem a caixa "decisao" sempre que registarem uma decisão
% importante — facilita a leitura e mostra o pensamento crítico
% da equipa.
% ---------------------------------------------------------------


\section{Visão geral}
\label{sec:arquitetura-geral}

% Visão de alto nível: três camadas, comunicação entre elas,
% protocolos. Acompanhar com o diagrama abaixo.

A solução adopta uma arquitetura em três camadas, separando
claramente as responsabilidades do cliente, do servidor e da
persistência. A Figura~\ref{fig:arquitetura-geral} apresenta a
visão de alto nível.

\begin{figure}[H]
\centering
\begin{tikzpicture}[
    node distance=1.6cm,
    bloco/.style={
      rectangle, rounded corners=3pt,
      draw=forest, line width=0.8pt,
      fill=boxbg,
      minimum width=3.4cm, minimum height=4cm,
      align=center, font=\small
    },
    seta/.style={
      <->, >=Stealth, line width=0.8pt, color=forest
    }
]

% Bloco BROWSER
\node[bloco] (browser) {%
  \begin{minipage}{3cm}
    \centering
    {\sffamily\bfseries\small\color{forest} BROWSER}\\[1pt]
    {\tiny\itshape\color{textMuted} Cliente}\\[6pt]
    {\scriptsize
      \texttt{HTML} / \texttt{CSS} / \texttt{JS}\\
      \texttt{localStorage}\\
      \texttt{sessionStorage}\\
      \texttt{IndexedDB}\\
      \texttt{Cache API}\\
      Service Worker
    }
  \end{minipage}
};

% Bloco API
\node[bloco, right=of browser] (api) {%
  \begin{minipage}{3cm}
    \centering
    {\sffamily\bfseries\small\color{forest} API NODE.JS}\\[1pt]
    {\tiny\itshape\color{textMuted} Servidor}\\[6pt]
    {\scriptsize
      Rotas REST\\
      Middleware JWT\\
      Validação Mongoose\\
      Motor de regras\\
      Logs auditoria\\
      OpenAPI
    }
  \end{minipage}
};

% Bloco BD
\node[bloco, right=of api] (db) {%
  \begin{minipage}{3cm}
    \centering
    {\sffamily\bfseries\small\color{forest} MONGODB}\\[1pt]
    {\tiny\itshape\color{textMuted} Persistência}\\[6pt]
    {\scriptsize
      Utilizador\\
      ErvaAromática\\
      PlanoCultivo\\
      LoteCultivo\\
      Tarefa, Medição\\
      Alerta, Auditoria
    }
  \end{minipage}
};

% Setas
\draw[seta] (browser) -- node[above, font=\tiny\sffamily, color=textMuted] {HTTPS · JSON} (api);
\draw[seta] (api) -- node[above, font=\tiny\sffamily, color=textMuted] {Mongoose} (db);

\end{tikzpicture}
\caption{Visão de alto nível da arquitetura em três camadas.}
\label{fig:arquitetura-geral}
\end{figure}

A comunicação entre o cliente e o servidor faz-se sobre HTTPS,
trocando mensagens em JSON. O servidor acede à base de dados
através do ODM \texttt{Mongoose}, que assegura a validação de
esquemas antes da persistência.


\section{Matriz de responsabilidades}
\label{sec:matriz-responsabilidades}

% Esta é a tabela mais importante do capítulo. Para cada
% funcionalidade, indica claramente o que fica em cada camada.
% Forçam-se assim respostas explícitas à pergunta "porquê aqui?"

A Tabela~\ref{tab:matriz} apresenta a distribuição de
responsabilidades por camada para as funcionalidades-chave do
sistema.

\begin{table}[H]
\centering
\caption{Distribuição de responsabilidades por camada.}
\label{tab:matriz}
\footnotesize
\renewcommand{\arraystretch}{1.4}
\begin{tabularx}{\linewidth}{l X X X}
\toprule
\textbf{Funcionalidade} & \textbf{Browser} & \textbf{API} & \textbf{Base de Dados} \\
\midrule
Autenticação JWT &
Envio de credenciais; armazenamento e renovação do token. &
Validação de credenciais; emissão e verificação do JWT. &
Persistência de utilizadores e \emph{hashes} de palavra-passe. \\

Registo de medição ambiental &
Captura via formulário; \emph{enqueue} em IndexedDB se \emph{offline}. &
Validação de regras; persistência; verificação de limites. &
Persistência da medição e dos alertas eventualmente gerados. \\

Geração de alertas &
Apresentação ao utilizador; UX de resolução. &
Avaliação dos limites do plano; criação dos alertas. &
Persistência do alerta com referência ao lote e ao plano. \\

Importação de CSV/Excel &
Pré-validação leve do ficheiro (formato, colunas). &
Validação completa, normalização e inserção em \emph{bulk}. &
Persistência das ervas aromáticas importadas. \\

Modo Manual / Automático &
UI de comutação e de listagem das ações sugeridas. &
Motor de regras; execução ou sugestão das ações. &
Registo da configuração do modo e do histórico de ações. \\

Sincronização \emph{offline} &
Fila de operações pendentes em IndexedDB; reenvio ao reconectar. &
Tratamento idempotente das operações; resolução de conflitos. &
Persistência idempotente; carimbos temporais. \\
\bottomrule
\end{tabularx}
\end{table}


\section{Camada de cliente (\emph{browser})}
\label{sec:camada-browser}

% Esta secção é o "documento de arquitetura cliente" exigido
% pelo enunciado. Tem de cobrir os 4 mecanismos de armazenamento.

A camada de cliente foi desenhada para suportar três cenários
distintos de operação: ligação estável, ligação intermitente e
ausência total de conectividade. Esta secção descreve os
mecanismos de armazenamento utilizados, a respetiva justificação,
política de invalidação e estratégia de sincronização.

\subsection{Mecanismos de armazenamento}

\begin{table}[H]
\centering
\caption{Mecanismos de armazenamento no \emph{browser} e respetivos propósitos.}
\label{tab:browser-storage}
\footnotesize
\renewcommand{\arraystretch}{1.4}
\begin{tabularx}{\linewidth}{l X X X}
\toprule
\textbf{Mecanismo} & \textbf{Dados guardados} & \textbf{Justificação} & \textbf{Expiração} \\
\midrule
\texttt{localStorage} &
Preferências da UI; identificador do último ecrã visitado; \emph{feature flags}. &
Pequenos pares chave/valor síncronos; persistem entre sessões. &
Manual; limpa no \emph{logout}. \\

\texttt{sessionStorage} &
Estado temporário de formulários multi-passo (\emph{wizards}). &
Volátil por aba; evita ``poluir'' o estado entre sessões. &
Fim da sessão da aba. \\

\texttt{IndexedDB} &
Histórico recente; cache de lotes e medições; fila de operações pendentes (\emph{offline}). &
Estrutura, queryable, assíncrono; permite operação degradada. &
TTL de [N] dias para o cache; a fila é purgada após sincronização bem-sucedida. \\

\texttt{Cache API} &
Recursos estáticos (HTML, CSS, JS, ícones); respostas \texttt{GET} idempotentes selecionadas. &
Gerida pelo Service Worker; reduz latência e suporta arranque \emph{offline}. &
Estratégia \emph{stale-while-revalidate}; invalidação por versão. \\
\bottomrule
\end{tabularx}
\end{table}


\subsection{Decisão sobre o armazenamento do JWT}
\label{sec:decisao-jwt}

% Aqui é OBRIGATÓRIO justificar com profundidade. O enunciado
% pede explicitamente esta decisão e a sua justificação de
% segurança. Substituir o exemplo abaixo pela decisão real da
% equipa.

\begin{decisao}[title=Decisão DA-01: armazenamento do token JWT]
\textbf{Decisão.} O token JWT é guardado em [\texttt{localStorage} / \texttt{sessionStorage} / \emph{cookie httpOnly}].

\textbf{Alternativas consideradas.}
\begin{itemize}
  \item \texttt{localStorage}: persistência longa, mas exposto a ataques XSS.
  \item \texttt{sessionStorage}: limita a janela de exposição mas obriga a re-autenticar a cada sessão.
  \item \emph{Cookie httpOnly + Secure + SameSite}: imune a XSS por JavaScript mas requer mitigação de CSRF.
\end{itemize}

\textbf{Justificação.} [Justificar com base no perfil de risco do
sistema: estufa em ambiente operacional, sessões longas, baixa
exposição a CSRF se a API não usa formulários tradicionais, etc.]

\textbf{Mitigações.} [Lista de medidas adotadas: \emph{Content Security
Policy} restritiva, \emph{escape} de saídas, sanitização de \emph{inputs},
expiração curta e renovação por \emph{refresh token}, etc.]
\end{decisao}


\subsection{Estratégia de sincronização}
\label{sec:sync}

% Como funciona o offline → online? Como se evitam duplicações?
% Como se resolvem conflitos?

A operação em modo degradado segue o padrão \emph{outbox}: as
operações realizadas \emph{offline} são enfileiradas em \texttt{IndexedDB}
e reenviadas ao servidor assim que a conectividade é restabelecida.
Para garantir a idempotência, cada operação carrega um identificador
único gerado no cliente (UUID) que o servidor reconhece e ignora em
caso de submissões repetidas.

\begin{aviso}
Em caso de conflito (e.g., duas equipas registam a mesma colheita
do mesmo lote em momentos diferentes), o sistema adota a estratégia
de [\emph{last-write-wins} / \emph{merge} / \emph{flag para revisão
humana}]. Esta decisão tem implicações operacionais e deve ser
discutida explicitamente.
\end{aviso}


\section{Camada de servidor (API)}
\label{sec:camada-api}

% Estrutura da API: organização das rotas, middleware,
% autenticação, motor de regras, geração de logs, geração da
% documentação OpenAPI.

\subsection{Organização das rotas}

A API expõe os recursos da aplicação seguindo as convenções REST.
A Tabela~\ref{tab:rotas} sintetiza os principais \emph{endpoints}
agrupados por recurso. A especificação completa encontra-se em
\texttt{openapi.yaml} (ver Anexo~\ref{anexo:openapi}).
Na Sprint 1 foi priorizada a especificação dos recursos de domínio,
ficando a implementação da autenticação para sprint posterior.

\begin{table}[H]
\centering
\caption{Resumo dos principais \emph{endpoints} da API.}
\label{tab:rotas}
\footnotesize
\renewcommand{\arraystretch}{1.3}
\begin{tabularx}{\linewidth}{l l X}
\toprule
\textbf{Método} & \textbf{Rota} & \textbf{Descrição} \\
\midrule
GET   & \texttt{/api/v1/users}                  & Listagem de utilizadores (Admin.). \\
POST  & \texttt{/api/v1/herbs/import}           & Importação em \emph{bulk} (CSV/Excel). \\
POST  & \texttt{/api/v1/cultivation-plans}      & Criação de plano de cultivo (\emph{regular}, \emph{emergência}, \emph{pontual}). \\
POST  & \texttt{/api/v1/measurements}           & Registo de medição ambiental com potencial geração de alerta. \\
POST  & \texttt{/api/v1/alerts/:id/ignore}      & Ignorar alerta com justificação obrigatória. \\
GET   & \texttt{/api/v1/reports/export}         & Exportação de relatório em CSV/Excel. \\
\bottomrule
\end{tabularx}
\end{table}


\subsection{Cadeia de \emph{middleware}}

Cada pedido atravessa uma cadeia de \emph{middleware} ordenada da
seguinte forma:

\begin{enumerate}
  \item \textbf{Logger HTTP}: regista método, rota, IP e \emph{user-agent}.
  \item \textbf{Validação JWT}: verifica assinatura e expiração; rejeita pedidos não autenticados em rotas protegidas.
  \item \textbf{Autorização por perfil}: confronta o perfil do utilizador (RF-01) com os requisitos da rota.
  \item \textbf{Validação de \emph{payload}}: valida o corpo do pedido contra o esquema Mongoose ou contra um \emph{schema} JSON.
  \item \textbf{Controlador da rota}: executa a lógica de negócio.
  \item \textbf{\emph{Middleware} de erros}: converte exceções em respostas HTTP coerentes com a OpenAPI.
\end{enumerate}


\subsection{Motor de regras e modo Automático}
\label{sec:motor-regras}

\begin{decisao}[title=Decisão DA-02: localização do motor de regras]
\textbf{Decisão.} O motor de regras corre na API, e não no cliente.

\textbf{Justificação.} Embora o cliente possa fornecer \emph{feedback}
otimista quando uma medição é registada \emph{offline}, a fonte da
verdade tem de estar no servidor: (i) os limites do plano podem
mudar; (ii) a coerência entre lotes exige uma visão global; (iii)
no modo \emph{Automático}, ações como ativar a rega ou ajustar a
ventilação têm efeitos físicos que não podem depender do estado
local do cliente.
\end{decisao}


\section{Modelo de dados}
\label{sec:modelo-dados}

% Apresentar o modelo conceptual com um diagrama (entidades e
% relações). Pode ser feito em TikZ ou desenhado num editor
% externo e incluído como imagem.

A Figura~\ref{fig:modelo-dados} apresenta o modelo conceptual
das entidades persistidas em MongoDB. Por simplicidade, são
omitidos atributos secundários.

% Substituir pelo diagrama real (e.g., draw.io exportado para PDF).
\begin{figure}[H]
\centering
\fbox{\parbox{0.9\linewidth}{\centering\vspace{2cm}
\textit{[Inserir aqui o diagrama de entidades e relações.]}\vspace{2cm}}}
\caption{Modelo conceptual de dados (entidades e relações).}
\label{fig:modelo-dados}
\end{figure}

\begin{nota}
Em MongoDB, cada documento pode ``embeber'' subdocumentos ou
referenciar outros documentos por \texttt{ObjectId}. As decisões
de \emph{embedding} versus \emph{referencing} influenciam o
desempenho das consultas e devem ser justificadas. Recomenda-se
uma subsecção dedicada onde, para cada par de entidades
relacionadas, se indique a opção tomada e o seu motivo.
\end{nota}


\section{Segurança}
\label{sec:seguranca}

% Reunir num só lugar todas as decisões de segurança.
% Útil para a equipa, útil para a avaliação.

As principais decisões de segurança adotadas são as seguintes:

\begin{itemize}
  \item As palavras-passe são armazenadas com \emph{hashing} adaptativo (\texttt{bcrypt}, fator de custo $\geq 10$).
  \item A autenticação usa JWT assinados (HS256 ou RS256) com \emph{payload} mínimo e expiração de [N] minutos.
  \item Todas as rotas, exceto as públicas, exigem JWT válido.
  \item A autorização é controlada por perfil; o servidor nunca confia em campos sensíveis enviados pelo cliente (e.g., o perfil do utilizador é sempre obtido do JWT, não do corpo do pedido).
  \item Os \emph{inputs} são validados contra esquemas; as saídas que possam conter texto introduzido pelo utilizador são \emph{escaped}.
  \item As ações sensíveis (criação de planos, autorização de planos pontuais, eliminação de utilizadores) ficam registadas em \texttt{LogAuditoria}.
\end{itemize}


\section{Resumo das decisões}
\label{sec:resumo-decisoes}

% Uma tabela síntese ajuda muito o avaliador a perceber o
% pensamento crítico da equipa. Listar as decisões principais
% com referência cruzada para a secção onde são justificadas.

\begin{table}[H]
\centering
\caption{Síntese das decisões de arquitetura.}
\label{tab:decisoes}
\renewcommand{\arraystretch}{1.3}
\begin{tabularx}{\linewidth}{l X l}
\toprule
\textbf{ID} & \textbf{Decisão} & \textbf{Secção} \\
\midrule
DA-01 & Armazenamento do JWT em [\dots]                              & \ref{sec:decisao-jwt} \\
DA-02 & Motor de regras na API                                       & \ref{sec:motor-regras} \\
DA-03 & Sincronização \emph{offline} via fila idempotente em IndexedDB & \ref{sec:sync} \\
DA-04 & [Adicionar conforme as decisões reais da equipa]             & --- \\
\bottomrule
\end{tabularx}
\end{table}
